% ============================================================================
% APPENDIX.  Eight sections, functional titles, coarse hierarchy.
% Protocol and extended results live TOGETHER under each experiment (F.1-F.6)
% rather than being split across a protocol appendix and a results appendix.
% Run-in labels are reserved for repeated technical schema; they are not used
% as rhetorical summaries or mini-results headings.
% ============================================================================

\section*{Overview of Appendix}
\addcontentsline{toc}{section}{Overview of Appendix}

This appendix provides the formal, implementation, and experimental details supporting the main text. \Cref{app:statespace} specifies the executable molecular state space, rewrite semantics, and canonical molecular-successor representation. \Cref{app:theory} gives the complete theoretical results underlying support preservation and finite-horizon control. \Cref{app:data,app:models,app:sampling} describe transition supervision, model training, and controlled sampling. \Cref{app:experiments} gives the experimental details and extended results for each experiment in the order it appears in \cref{sec:experiments}, \cref{app:alignment} documents alignment with external benchmark protocols, and \cref{app:repro} records the software environment, randomization scheme, and compute accounting.

\section{Molecular state space and executable edits}
\label{app:statespace}

The support-level claims of \method{} are defined with respect to an explicit molecular state space rather than post-hoc sanitization. This section specifies that state space and the rewrite contracts that induce its transition graph. The \emph{semantic state} is the active molecular graph on which connectivity, valence, ring perception, and canonical identity are evaluated, and we distinguish it from any padded tensor representation used by the implementation. We then define the current operator family, its exact legality checks, and the pushforward from executable rewrite marks to canonical molecular successors. These definitions determine what transitions are possible before either $R_\theta$ or a task-specific controller assigns probability to them.

\subsection{Molecular state space}

States carry at most \XXX{} heavy atoms over \XXX{} neutral (element, valence) classes drawn from a declared organic element set, a vocabulary admitting \XXX{} of the benchmark lead molecules.
Formal charge is an input but not an output dimension, so charge is carried rather than designed; stereochemistry, isotopes and radicals are not represented.

Cardinality is a property of the active molecular graph rather than of a fixed tensor allocation. Writing $\X_n$ for the states with $n$ heavy atoms, the state space is the disjoint union of the strata, and a birth or death moves a trajectory between strata by creating or destroying a typed graph entity that participates in valence, connectivity, and ring perception. This is not the activation or deactivation of a padded coordinate, and the distinction is what makes the process trans-dimensional in the sense used throughout the paper.

\subsection{Rewrite operators}
\label{app:operators}

The main text depends only on the legal edit set \eqref{eq:fiber} and never on which families populate it; this section records the families used, so that the edit set is reproducible. We use paper-facing operator names for executable rewrite semantics and model-family names when reporting learned-family statistics; \cref{tab:operators} gives the fixed mapping between them.

\begin{table}[h]
\centering
\small
\caption{The rewrite families, by action on the strata. Each family is a typed graph rewrite with structural and chemical application conditions; the executor follows a propose--validate--commit discipline and exposes no transition whose successor fails a guard. Names here are the paper-facing operator names. The implementation additionally carries model-family identifiers, which appear in per-family results: $\op{cycle\_close}$ is recorded as \texttt{cycle\_insert}, $\op{cycle\_open}$ as \texttt{cycle\_attach}, and $\op{ring\_aromaticity\_restate}$ as \texttt{ring\_system\_restate}. The remaining five names coincide in both registers. Per-family results in \cref{app:exp-rtheta} are reported under the model-family identifiers.}
\label{tab:operators}
\begin{tabularx}{\textwidth}{@{}llY@{}}
\toprule
Family & Cardinality & Structural / chemical condition\\
\midrule
$\op{atom\_insert}$ & birth $\X_n\!\to\!\X_{n+1}$ & empty capacity; a null source or exactly one occupied anchor; allowed atom state; for a bonded birth, allowed bond order and feasible endpoint valence\\
$\op{atom\_delete}$ & death $\X_n\!\to\!\X_{n-1}$ & removable atom with connected remainder; neighbour hydrogen compensation; successor sanitizes\\
$\op{atom\_restate}$ & same & occupied site; joint element, valence-class and implicit-hydrogen state allowed; successor differs from the source\\
$\op{bond\_reorder}$ & same & existing edge; joint bond-order and hydrogen update feasible\\
$\op{bond\_reroute}$ & same (reroute) & the cut edge is a bridge; the relocated fragment is a pendant subtree; the new endpoint lies outside it; cut and attach commit in one transaction so no disconnected state is ever exposed\\
$\op{cycle\_close}$ & same (cyclize) & absent edge between two occupied atoms with valence headroom; admitted insertion raises graph cycle rank by one\\
$\op{cycle\_open}$ & same (decyclize) & existing edge that is not a bridge; endpoint hydrogen compensation; successor sanitizes\\
$\op{ring\_aromaticity\_restate}$ & same & a matched cyclic system; a coordinated set of bond-order changes validated by the executor as one transaction\\
\midrule
whole-ring growth, whole-ring deletion & multi-atom macros & \textbf{disabled}; a macro must earn its place by a measured path-efficiency gain, never by assumption\\
\bottomrule
\end{tabularx}
\end{table}

$\op{cycle\_close}$ inserts one bond between two nonbonded atoms and, because every non-null state is connected, raises graph cycle rank by one; $\op{cycle\_open}$ deletes one non-bridge cycle edge and lowers it by one.
At the untyped graph level every connected graph is a spanning tree plus its non-tree edges, so inserting those edges one at a time realizes the abstract cyclic topology while preserving connectivity.
This is a statement about graphs.
It does not prove that the molecular grammar admits a valence-valid path to every chemically labelled ring system, nor that such a path fits a finite edit budget; both are measured questions and neither is assumed anywhere in the main text.

\subsection{Canonical molecular successors}

Distinct executable marks can produce the same molecule, so the law the controller acts on is the pushforward of the mark law through the executor, aggregated at canonical molecular identity as in \eqref{eq:sucfiber}--\eqref{eq:rtheta}. Aggregation is what makes every downstream quantity independent of how many ways the implementation can name a molecule; \cref{app:quotient} establishes the precise invariance this buys and its exact scope.

\section{Theoretical results}
\label{app:theory}

Results are ordered as a reader needs them: state-space closure first, then the canonical-successor quotient and its invariance, then the finite-horizon control results that depend on both. Each statement is repeated here with its assumptions so that the appendix stands alone.

\subsection{Executable-support closure}
\label{app:proofs}

We restate each result with its assumptions before proving it, so that this appendix can be read independently of the main text.

\begin{proposition}[Closure of executable support; restatement of \cref{prop:closure}]
Membership in $\X$ requires exactly four conditions: node and edge attributes lie in the declared vocabularies, every atom satisfies its declared valence constraints, the graph is connected, and its heavy-atom count lies in $\{1,\dots,n_{\max}\}$. Suppose each admitted edit keeps every created or modified attribute within the declared vocabularies, preserves the declared valence constraints at every atom it modifies, leaves untouched atoms unchanged, respects the heavy-atom bound $n_{\max}$, and is committed only when its successor is connected. Then $\T(x,a)\in\X$ for every $x\in\X$ and $a\in\A(x)$, and consequently every state reached by any stochastic process supported on $\{\A(x)\}_{x\in\X}$ remains in the declared molecular state space.
\end{proposition}

\begin{proof}
Induct on committed events. The initial state lies in $\X$. If $x_n\in\X$, the process has support only on $a\in\A(x_n)$, and by \eqref{eq:fiber} every such mark executes to a state that passed the executor's guards. Those guards are exactly the four membership conditions above, checked before the edit is committed, so the successor lies in $\X$. The same induction applies verbatim to the controlled process of \eqref{eq:central}, because the controlled off-diagonal mass is a pointwise product of a nonnegative factor with a base probability and therefore vanishes wherever the base probability vanishes: no controller creates support.
\end{proof}

\subsection{Canonical successor representation}
\label{app:quotient}

\begin{definition}[Successor-preserving, rate-preserving re-encoding]
\label{def:reencoding}
Two mark systems on the same state space are so related if for every $x$ there is a surjection $\varphi_x$ from the marks of one onto the marks of the other with matching executed successors, so that marks may be split or merged only \emph{within} a successor fiber and never across fibers, and if the aggregate rate of every successor fiber agrees.
\end{definition}

\begin{proposition}[Encoding invariance, and its exact scope]
\label{prop:encoding}
Under \cref{def:reencoding} and a common state-level productive-commit predicate:
\emph{(i)} the state-level generators coincide, hence so do the total hazards, the molecular laws, the reference edit law $R_\theta$ of \eqref{eq:rtheta} wherever defined, and every time marginal;
\emph{(ii)} for any controller that is \emph{fiber-constant}, meaning a function of $(x,y)$ rather than of the mark, the controlled kernels coincide;
\emph{(iii)} hard state predicates and the transform of \cref{thm:doob} are fiber-constant by construction, so the values $h_b$ and the controlled kernels are identical under any such re-encoding.
\end{proposition}

\begin{proof}[Proof]
\emph{(i)} The fibers correspond under $\varphi_x$ and partition the applicable set, so fiber sums agree successor by successor; summing over successors gives equality of total hazards without further hypothesis, which answers the natural worry that a refinement, which renormalizes the mark law over a larger set, could perturb the hazard. It cannot, because the hazard is the sum of the aggregates. Equality of generators plus uniqueness of the finite-state forward solution gives equality of all marginals, and dividing the aggregate by the hazard gives equality of $R_\theta$ on every row with positive mass.
\emph{(ii)} A fiber-constant controller is a function $\Psi(x,y)$, so it multiplies two equal aggregates and yields equal controlled rates; the diagonal is determined by the row.
\emph{(iii)} $h_b$ is built by backward recursion from kernels that are invariant by \emph{(i)}, so $h_b$ is invariant by induction, and the ratio $h_{b-1}(y)/h_b(x)$ is fiber-constant.
\end{proof}

\begin{remark}[Scope of encoding invariance]
\Cref{prop:encoding} applies to successor-preserving refinements and to controllers defined on molecular successors. It does not imply invariance of arbitrary mark-level control.
\end{remark}

Merging marks \emph{across} successor fibers is excluded by \cref{def:reencoding} and is ill-posed rather than merely lossy: it would stop the executor being a function of the merged mark.
And mark-level controllers are not invariant in general.
If a successor is reachable by four symmetric aliases of rate $1/4$ under one encoding and by a single mark of rate $1$ under another, a mark-level power tilt $\Psi\propto q(a)^{\beta-1}$ assigns it controlled mass $4\cdot(1/4)^\beta$ versus $1$, which differ for every $\beta\ne1$.
Top-$k$ and nucleus truncation over marks, and per-family multipliers, fail for the same reason; a family multiplier is invariant exactly when it is constant across every alias of every successor fiber it touches.
Every controller in this paper is written over molecules.
The proposition concerns the pushed-forward process, not the training representation: mark-level supervision is generally not invariant to within-fiber refinement, which is why teacher supervision aggregates the successor fiber.

\subsection{Relation to generator matching}
\label{app:gm}

\Cref{sec:rtheta} trains $R_\theta$ by the productive-successor likelihood \eqref{eq:refloss}. This section records the exact relationship between that objective and generator matching \citep{holderrieth2025generator}, and the two respects in which they differ.

Write the productive rate at a state as $Q^{+}(y\mid x,t)=\Lambda(x,t)\,P(y\mid x,t)$, separating a total productive exit rate from a normalized law over canonical molecular successors. For the Bregman divergence generated by $\varphi(u)=u\log u-u$, a direct computation gives, for each $(x,t)$,
\begin{equation}
D_\varphi\!\left(Q^{+\star},Q^{+}_\theta\right)
=\underbrace{\Lambda^{\star}\,\KL\!\left(P^{\star}\,\middle\|\,R_\theta\right)}_{\text{conditional-successor block}}
\;+\;\underbrace{\Bigl(\Lambda^{\star}\log\tfrac{\Lambda^{\star}}{\Lambda_\theta}-\Lambda^{\star}+\Lambda_\theta\Bigr)}_{\text{productive-hazard block}},
\label{eq:gmsplit}
\end{equation}
where the second term is the same scalar divergence applied to the exit rates. The two blocks are variationally separable: the first depends on $\theta$ only through $R_\theta$, the second only through $\Lambda_\theta$. Minimizing over the successor law therefore leaves the hazard block untouched, and conversely.

\method{} optimizes an empirical version of the conditional-successor block, and differs from the generator-matching population objective in exactly two respects. First, the hazard block is never optimized. The compiled supervision carries no rate magnitudes: every stored transition is registered with unit teacher rate, which enters training only as a nonterminality indicator, so no target for $\Lambda_\theta$ exists in the corpus. Under that unit-rate convention $\Lambda^{\star}$ is constant and the weight it contributes to the first term is absorbed into a positive constant independent of $\theta$. The reported results never read an exit-rate magnitude; fixed-budget editing consumes only the embedded successor kernel, on which a constant rescaling of $Q^{+}$ has no effect.

Second, the expectation is taken under a different measure. Generator matching integrates the per-state divergence against the state--time marginal of the conditional-process mixture. \method{} instead draws training rows under a declared editing measure that is weighted across operator families by design, so that a family's exposure reflects the reference law the model is meant to estimate rather than the incidental composition of the compiled corpus; some strata are assigned zero draw. No importance correction is applied, and on zero-draw strata none could be. Consequently \eqref{eq:gmsplit} holds pointwise in $(x,t)$, and \eqref{eq:refloss} is the conditional-successor block of that decomposition, but the training functional is not the generator-matching population objective. We therefore describe the relationship as a decomposition that identifies which factor $R_\theta$ estimates, not as an objective identity.

\subsection{Finite-horizon control}

\begin{theorem}[Finite-horizon molecular control; restatement of \cref{thm:doob}]
Let $R_\theta$ define the reference path measure over $K$ committed transitions from an initial molecule $x_0$, with trajectories terminating before $K$ assigned zero continuation mass, let $g_z\ge0$ be a finite terminal desirability, and let $h_b$ be given by the backward recursion \eqref{eq:backward}. Wherever those backward values are positive, the controlled kernels $R^{z,b}_\theta$ of \eqref{eq:doob} are normalized, introduce no transitions outside the support of $R_\theta$, and after $K$ controlled edits realize the terminal law \eqref{eq:tilt}.
\end{theorem}

\begin{proof}
\emph{(i)} $\sum_y R_{\theta,b}(x,y)h_{b-1}(y)=h_b(x)$ by \eqref{eq:backward}, so the ratios in \eqref{eq:doob} sum to one.
\emph{(ii)} Multiplication by $h_{b-1}(y)$ cannot make a zero base probability positive; the supports are equal at a row exactly when no base-supported successor has zero value.
\emph{(iii)} For a path $x_0,\dots,x_K$ of positive base probability, the controlled path probability is the base path probability times $\prod_{k}h_{K-k-1}(x_{k+1})/h_{K-k}(x_k)$, and the interior factors telescope, leaving $h_0(x_K)/h_K(x_0)=g_z(x_K)/h_K(x_0)$. Every path to a fixed endpoint is therefore reweighted by the same factor, and summing over paths gives \eqref{eq:tilt}. The normalizer is the full-budget value at the start.
Finally, for the continuation and horizon-specificity claims stated after \cref{thm:doob}, apply the same telescoping to the chain started at $x$ with horizon $b'$ and desirability $g_{z'}$. The marginal at $j<K$ is an $h_{K-j}$-tilt rather than a $g_z$-tilt, which is why early stopping does not inherit \eqref{eq:tilt} and why continuation is not the same law as using $g_{z'}$ from the first edit.
\end{proof}

\subsection{Retargeting and pathwise support}

Retargeting is \cref{thm:doob} applied from a realized state. Instantiating the theorem at $(x_\tau,b)$ with desirability $g_{z'}$ gives exactly the $g_{z'}$-tilt of the reference law started at $x_\tau$ with budget $b$, which is in general not the law obtained by specifying $z'$ at the source lead; the transform applies because $x_\tau$ is itself a molecule by \cref{prop:closure}, so $g_{z'}$ is decidable on it.

A rule-closed predicate $C$ restricts the legal edit set to $\A_C(x)$ as in \eqref{eq:constrained}. Because the restricted edit set is again a set of executable marks into $\X$, the closure argument above applies unchanged to the restricted process, \cref{thm:doob} holds on it, and $C$ holds at every committed molecule by induction rather than by filtering at the end. Cumulative requirements are not predicates on the current molecule; augmenting the state with a running statistic $\xi_{k+1}=\Phi(\xi_k,X_{k+1})$ restores the Markov property on the augmented chain, and the same results then apply to $\widetilde X_k=(X_k,\xi_k)$.


\section{Data and transition supervision}
\label{app:data}

A reader should be able to answer, from this section alone, what evidence taught $R_\theta$ to prefer one legal transition over another. That question has two parts: which molecules the transitions were drawn from, and how each transition was constructed. The second part is not uniform across the edit space, and it bounds how the strongest per-family results in \cref{app:exp-rtheta} may be read.

\subsection{Molecular data and preprocessing}

The editing corpus was derived from a fixed $500{,}000$-molecule GuacaMol source pool by compiling executable local transitions, real-endpoint analogue relations, and reversible synthetic walks. \method{} was not trained as an unconditional distribution model over that pool: the learned object is $R_\theta(y\mid x)$, a transition law over source--successor edits, not a density $p_\theta(x)$ over molecules.

The source pool is identified by the recorded dataset namespace and its SHA-256, and independently corroborated by the molecular statistics of the frozen training sources. Measured on $4{,}000$ randomly sampled training sources, the element distribution is exactly the declared organic vocabulary, with carbon at $0.748$, nitrogen and oxygen at $0.105$ each, sulfur at $0.015$, fluorine at $0.012$, chlorine at $0.009$, bromine at $0.003$, iodine and phosphorus at $0.002$, and boron at $0.001$. The median heavy-atom count is $26$, median molecular weight $371.6$ and ninetieth percentile $499.3$, median ring count $3$, median aromatic-ring count $2$, and only $1.6\%$ acyclic. Because the training receipts do not bind the original raw asset end to end, we treat this provenance link as strongly supported rather than cryptographically complete; the frozen compiled corpus itself is exactly identified by its dataset, library, packed-store and reserve digests. The benchmark sources of \cref{app:bench-qed} are drawn from a different pool.

\subsection{Transition construction}

Transitions enter the corpus through several lanes, which are near family-pure with lane-homogeneous chunks. The frozen structural gate records $1{,}803{,}032$ training transitions passing with zero violations. A metadata scan of every training chunk attributes \texttt{cycle\_insert} ($122{,}183$ transitions across $102$ chunks), \texttt{cycle\_attach} ($84{,}940$ across $102$) and \texttt{ring\_system\_restate} ($16{,}380$ across $29$) entirely to the reversible synthetic-walk lane, with zero transitions in any of the four data-backed lanes; in the pinned distribution-matched training subset that lane contributes $23{,}100$ of $145{,}756$ transitions, or $15.8\%$, against $84.2\%$ data-backed. By contrast \op{bond\_reroute} and \op{atom\_restate} each appear in two lanes and therefore do have a real-chemistry source. The analogue lane draws on a frozen molecular-matched-pair pool of $344{,}153$ records bound by a recorded contract hash, whose own upstream molecular source is not fully documented.

The consequence is a scope limit rather than a defect, and it is applied throughout: evidence from the ring and cycle families supports the claim that the process recovers legal ring edits under synthetic perturbation, and does not support a claim of data-backed ring editing or scaffold hopping. Lifting that scope requires a data-backed ring-transformation lane, which the current corpus does not contain.

\subsection{Supervision provenance and data splits}

Every result in \cref{sec:experiments} is bound to artifacts frozen before it. The reference edit law was never retrained for any experiment in this paper, and the controller head evaluated on the prospective panel was verified byte-identical by hash before and after an unrelated training session. The three source panels (development, prospective validation, and the official cohort) were verified pairwise disjoint, with the official cohort and the development panel excluded from the eligible pool before the training corpus was sampled. Role precedence is final test, then controller validation, then validation, then train, so a molecule appearing in two raw libraries is excluded from the lower-precedence role; the one such case is excluded from training and appears in neither the final training set nor the matched reserve.

Both evaluation panels are exactly disjoint and near-duplicate disjoint from the training source universe of $96{,}094$ molecules. The matched reserve is cluster-disjoint by construction under an ECFP4 Tanimoto threshold of $0.95$ with formula and scaffold blocking, and the external cohort shows zero exact overlap under a scaffold-matched comparison. The reserve's higher scaffold sharing is by design: it was carved from the same population, which is what makes it population-matched. The evaluation ladder and its no-rescue rules were recorded before the data: the validation panel is opened once, whatever it returns is banked, and no hyperparameter is re-chosen afterwards.

A pre-freeze mechanical smoke test executed one official benchmark source before the stop condition took effect. Its output was never inspected, is excluded from every analysis, and is retained only as an audit artifact; we therefore describe the official cohort as inferentially unopened rather than literally never executed. The sealed final-test role, comprising $33{,}683$ source records, has not been leakage-audited, that audit is a precondition for opening it, and no result in this paper draws on it.

\section{Model and training details}
\label{app:models}

Three learned objects appear in this paper and are trained separately: the goal-independent reference law $R_\theta$, a property model $\widehat F_\psi$ that defines molecular desirability, and goal- and budget-conditioned future-value models. Only the future-value models are trained against a downstream design objective; $R_\theta$ never sees one.

\subsection{Reference edit law}

$R_\theta$ is fitted once, frozen, and never retrained for any experiment reported here. Its identity is pinned by the sampling-law hash recorded in the corpus binding artifact and by the dataset, library, packed-store and reserve digests alongside it. The conditional edit law is normalized over the legal edit set as in \eqref{eq:marklaw} and is hierarchical as in \eqref{eq:hierarchical}: a shared molecular encoder produces a state representation of width $256$, a family gate scores which rewrite family fires, and family-specific scorers rank legal operands and payloads within it. Normalization is restricted to candidates the executor has admitted, so the model never places mass outside $\A(x)$. The encoder reads the current molecular graph and a progress variable, and does not read the source lead, the objective, the remaining budget, or any task constraint.

The reference law is trained by maximum likelihood of the next committed molecular successor under the declared editing training measure. For a productive teacher transition $x_i\rightarrow y^{\star}_i$,
\begin{equation*}
\mathcal{L}_{\mathrm{ref}}(\theta)=-\frac{\sum_i w_i\log R_{\theta,t_i}(y^{\star}_i\mid x_i)}{\sum_i w_i},
\end{equation*}
as in \eqref{eq:refloss}, where $R_{\theta,t}$ is the canonical successor kernel of \cref{sec:canonical} and terminal rows are excluded. For the frozen checkpoint used throughout this paper, teacher rates and importance weights are unit-valued; the corpus sampler instead defines the reference measure through its registered family-weighted row distribution. The total event intensity is not optimized by this objective, and fixed-budget \method{} therefore does not claim calibrated continuous-time holding rates.

% ===== GAP: NUMBER NOT IN VERIFIED SOURCE SET =====  [U4] R_theta training constants
% Encoder depth and width, message-passing steps, optimizer, schedule, batch
% size, total steps and checkpoint selection are recorded in the training
% configuration and run receipts on the artifact volume.  Design constants, not
% measurements, but outside this draft's verified source set.
\emph{Training.} Encoder depth \XXX{}, message-passing steps \XXX{}, optimizer and schedule \XXX{}, batch size \XXX{}, total optimization steps \XXX{}, checkpoint-selection criterion \XXX{}.

\subsection{Property models}

Molecular desirability is defined through a frozen property model $\widehat F_\psi$ rather than through oracle calls at control time. For the two-objective experiment it is fitted to a matched budget of $10{,}000$ random-ZINC property labels, the same budget the comparator receives, and is shared identically by the immediate and future-aware arms so that any difference between them is attributable to control rather than to the landscape. Both arms therefore inherit the same prediction error and the same extrapolation behaviour outside the label range.

% ===== GAP: NUMBER NOT IN VERIFIED SOURCE SET =====  [U5] property-model constants
\emph{Architecture and training.} Architecture \XXX{}, featurization \XXX{}, held-out predictive performance \XXX{}.

\subsection{{QED} future-value model}
\label{app:controller}

The region-goal controller estimates finite-budget reachability of a goal region and is used for the similarity-constrained editing experiments. The values below are design constants fixed in frozen preregistration documents before the corresponding data existed.

\emph{Inputs.} The controller conditions on the region's own coordinates, the signed margins of the current molecule to each of them, the remaining budget as a one-hot index, and a graph embedding of the current molecule together with the source lead and their difference and product, produced by the frozen reference-law encoder used unchanged with gradients never flowing into it. The feature width is $4\times256+2+2+2+(24+1)=1055$: four embedding blocks, two state properties, two region thresholds, two signed margins, and a $25$-slot budget one-hot. The budget ceiling of $24$ is frozen because the released checkpoint was trained against a $25$-slot one-hot; a longer horizon is expressed by passing a larger ceiling to the feature builder rather than by editing the constant, which would change the input width and make the checkpoint unloadable.

\emph{Architecture and objective.} The head is a multilayer perceptron $1055\rightarrow512\rightarrow256\rightarrow128\rightarrow1$ with rectified-linear activations and dropout $0.1$ after the first two hidden layers. It emits a logit; the sigmoid is applied outside the module, so the multiplicative factor in \eqref{eq:central} lies in $[0,1]$ by construction. The loss is binary cross-entropy against the Monte-Carlo hitting label plus a Bellman consistency term at weight $0.3$ on observed transitions, whose target is the sigmoid of the head at the successor. Optimization uses Adam at learning rate $10^{-3}$ with weight decay $10^{-5}$ for at most $40$ epochs with early stopping at patience $5$, and the best-validation checkpoint is selected rather than the last epoch.

\emph{Goal conditioning.} Conditioning is continuous in the goal, so a threshold pair never seen in training is a point in the same input space rather than an unseen class label. The registered grid is the product of five QED thresholds $\{0.75,0.80,0.85,0.90,0.95\}$ and four similarity floors $\{0.30,0.40,0.50,0.60\}$, giving $20$ jointly trained goal regions; the benchmark's region $(0.90,0.40)$ is one member and is given no special weight. The boundary condition $h_\phi\equiv1$ on the region is enforced at inference, and in-region states are excluded from training rather than supervised toward it.

\emph{Splitting.} The validation fraction is $0.15$, held out by source molecule rather than by training example, because a trajectory rolled out of one lead yields many near-duplicate states and holding out a random fraction of examples would report memorization as generalization.

\subsection{Multi-objective future-value model}
\label{app:mo-controller}

The multi-objective controller estimates a different quantity from the region model above and is documented separately for that reason. Its target is the expected preference-weighted terminal desirability under the frozen reference process, $h_b(x,\lambda)=\E_{R_\theta}[g_\lambda(X_b)\mid X_0=x]$, with $g_\lambda$ the Chebyshev desirability of \eqref{eq:glambda}; it is a soft expectation rather than a hitting probability, so the binary construction of the region model does not carry over.

For numerical conditioning, the implementation evaluates the equivalent centered form $\widetilde g_\lambda(x)=\exp\bigl(-(L_\lambda(x)-L^{\star}_\lambda)/\tau_\lambda\bigr)$, where $L^{\star}_\lambda$ is a per-preference constant with respect to $x$. This multiplies all desirabilities for a fixed preference by the same positive constant, so the factor cancels in the normalization of \eqref{eq:central} and leaves every controlled transition unchanged.

Supervision is Monte Carlo: $M$ independent continuations of the frozen reference law from a state give the estimator of \eqref{eq:mctarget}. The variance of that estimator grows with the remaining budget while the spread of true values across candidate successors shrinks, so target reliability degrades with horizon. \Cref{app:exp-pareto} reports that degradation and the rollout counts it implies.

% ===== GAP: NOT YET RUN =====  [G3] multi-objective future-value model
% The supported-horizon gate, the rollout count admitted at each horizon, and
% the trained multi-objective head do not yet exist.  TO CLOSE: fix the minimum
% target-reliability criterion, determine the rollout count meeting it at each
% candidate horizon, and train only on admitted horizons.
Because supervision quality is horizon-dependent, learned future-value control is admitted only at horizons whose target reliability meets a criterion fixed in advance of any optimization result. The admitted horizon $b_{\mathrm{FA}}$ and the corresponding rollout count are \XXX{}, and the deployed score is $h_\phi(y,\lambda,b-1)$ for $b\le b_{\mathrm{FA}}$ and $g_\lambda(y)$ otherwise.

\subsection{Training-pipeline schematic}

% APPENDIX FIGURE --- training / data-flow schematic.  Deliberately NOT promoted
% to the main text: Figure 1 plus the Method prose already carry this, and a
% second schematic before any evidence reads as schematic -> schematic -> results.
\begin{figure}[h]
\centering
\fbox{\begin{minipage}{0.95\linewidth}\centering\footnotesize
{\color{red}\textbf{[Schematic to be drawn. No data plotted.]}}\\[3pt]
\textbf{1.} pairs $\to$ certified edit programs $\to$ successor-likelihood training $\to R_\theta$ \emph{(frozen)}\\
\textbf{2.} frozen-$R_\theta$ rollouts $+\,z\to$ future-value labels $+$ Bellman $\to h_\phi$\\
\textbf{3.} $\widehat F_\psi$ read at the successor \emph{vs.} $\widehat F_\psi$ propagated through $R_\theta$
\end{minipage}}
\caption{\small\textbf{What is trained, and what each trained object is asked to know.} The reference edit law is fitted from certified programs of executable edits and frozen thereafter, reading neither the source lead nor the objective; the goal-conditioned value is fitted from rollouts of that frozen law, with the budget axis explicit and validation held out by source molecule; the third row is the immediate-versus-future-aware contrast. Nothing here is a result.}
\label{fig:trained}
\end{figure}


\section{Sampling and control}
\label{app:sampling}

Two realizations of the controlled law \eqref{eq:central} are used. One is exact given the learned value; the other is an approximation whose error is in particle allocation rather than in the definition of success. Neither is a separate contribution.

\subsection{Exact finite-horizon control}

Where the set of molecular states reachable within the remaining budget can be enumerated, the backward value of \eqref{eq:backward} is computed by exact dynamic programming: initialize $h_0=g_z$ on the enumerated terminal layer and sweep backwards, at each budget taking the expectation of the next layer's value under $R_\theta$ restricted to the legal edit set. The controlled kernel of \eqref{eq:doob} follows by division and \cref{thm:doob} applies without approximation. This is the only setting in which the word \emph{exact} is used of a deployed computation.

\subsection{Molecular-scale sampling}

At scale $h_\phi$ replaces the exact value. Because it is a state-level function taking values in $[0,1]$, the controlled row admits a rejection realization: propose $Y\sim R_\theta(\cdot\mid x)$, execute, canonicalize, reject self-transitions, and accept with probability $h_\phi(Y,z,b-1)$. Conditioned on acceptance the move has exactly the $R_\theta\cdot h_\phi$ law, and aliased successors aggregate correctly because the acceptance probability depends only on the resulting molecule. The proposal cap is finite, however, so when every proposal at a state is rejected the trajectory ends; the capped procedure is therefore not an exact realization of the controlled kernel and is not described as one. Cap pressure is recorded as a diagnostic with a preregistered escalation to the sequential realization below, which changes inference only.

When the desired region carries little mass under the reference process, control is realized as a twisted Feynman--Kac particle system. Every particle proposal comes from the frozen reference law alone, and $h_\phi$ enters only through the incremental weight \eqref{eq:fk}. Using the value to both select and weight would double-count control and sample a different process. The weight product telescopes along a trajectory to $h_0(x_H)/h_H(x_0)$, and the terminal boundary $h_0(x)=\one[x\in\Bset_z]$ is exact, so a learned twist changes particle allocation without changing the terminal success criterion. The operating point is frozen and no particle-count sweep is permitted: $N=32$ particles, horizon $24$, absorbing stop on entry to the goal region, and systematic resampling whenever the effective sample size $\mathrm{ESS}(w)=(\sum_i w_i^2)^{-1}$ on normalized weights falls strictly below $N/2$. Resampling is systematic rather than multinomial, a single uniform draw positioning the whole comb. Per-transition values, per-step resampling decisions and resampling indices are recorded so the mechanical checks run offline against the artifact.

Candidate accounting follows one rule: one run returns one molecule. A source's twenty candidates are twenty independent runs, each contributing the single molecule drawn from its terminal normalized particle measure. The particles within a run are inference state and are never counted as candidates, and intermediate states are never retrospectively harvested. A run reaching the budget without entering the region is recorded under a frozen extinction status rather than replaced. For a region goal the controller stops at the first committed molecule inside the region, an online stopping time evaluated on the realized state.

\begin{algorithm}[h]
\caption{Controlled molecular sampling, sequential realization}
\begin{algorithmic}[1]
\STATE \textbf{input:} source $x_0$; goal $z$; budget $K$; frozen $R_\theta$ and $h_\phi$; particles $N$; optional rule-closed predicate $C$
\STATE initialize $x^i \leftarrow x_0$ and $w^i \leftarrow 1/N$ for $i=1,\dots,N$; \ $b \leftarrow K$
\WHILE{$b > 0$ and some particle is unabsorbed}
  \FOR{each unabsorbed particle $i$}
    \STATE enumerate $\A(x^i)$; if $C$ is given, restrict to $\A_C(x^i)$
    \STATE push forward through the executor and canonicalize to $\Sset(x^i)$
    \STATE draw $y^i \sim R_\theta(\cdot \mid x^i)$ \hfill \COMMENT{proposal from the reference law alone}
    \STATE $w^i \leftarrow w^i \cdot h_\phi(y^i,z,b-1) \,/\, h_\phi(x^i,z,b)$ \hfill \COMMENT{the value enters only here}
    \STATE $x^i \leftarrow y^i$; \ if $x^i \in \Bset_z$ then absorb particle $i$
  \ENDFOR
  \STATE normalize $w$
  \IF{$\mathrm{ESS}(w) < N/2$}
    \STATE systematically resample the particles and reset $w^i \leftarrow 1/N$
  \ENDIF
  \STATE $b \leftarrow b-1$
\ENDWHILE
\STATE \textbf{return} one molecule drawn from the terminal normalized particle measure, or the frozen extinction status if no particle was absorbed
\end{algorithmic}
\end{algorithm}

\subsection{Multi-objective control}

Under a preference $\lambda$ the immediate arm weights successors by $g_\lambda(y)$ and the future-aware arm by $h_\phi(y,\lambda,b-1)$, with the reference law, legal edit set, property model, sources, preferences and budget identical between them. Where the admitted horizon $b_{\mathrm{FA}}$ is shorter than the optimization campaign, a campaign consists of repeated controlled trajectories of length at most $b_{\mathrm{FA}}$ launched from the frozen initialization bank and its archive until the common evaluation budget is exhausted. The campaign horizon and the horizon over which learned foresight is claimed are therefore distinct quantities.

\subsection{Retargeting and pathwise constraints}

Retargeting replaces the goal context at a realized state and continues with the remaining budget by re-evaluating \eqref{eq:central} at that state; no parameter of the reference process changes and the prefix is retained. A rule-closed constraint is imposed by filtering the legal edit set to $\A_C(x)$ of \eqref{eq:constrained} before the transition is sampled, so no infeasible successor is proposed, weighted, or committed.


\section{Experimental details}
\label{app:experiments}

Each subsection below gives the setup, arms and statistics for one experiment, together with the extended results omitted from the main text, in the order the experiments appear in \cref{sec:experiments}.

\subsection{Reference-law evaluation}
\label{app:exp-rtheta}

We evaluate the frozen reference law on held-out transitions from the matched reserve by measuring the probability assigned to the realized canonical successor. The uniform, empirical-family and learned laws are evaluated on the identical executable support and canonical successor quotient; only their probability assignments differ. The empirical-family law uses corpus-level edit-family frequencies renormalized over the families available at the current state and distributes each family's mass uniformly across its available successors. Two intermediate arms decompose the comparison: an empirical family with learned identity, which keeps corpus family frequencies and learns only which successor within the chosen family is realized, and a learned family with uniform identity, which learns only which family fires. Together they separate choosing which edit to make from choosing where to apply it.

The primary quantity is the mean negative log-probability of the realized successor, with the scored transition as the unit and paired differences by construction, since every law scores the same transition. Cells below a preregistered minimum size of $20$ are not reported separately; no scored transition was skipped and no partition failed to compile. Performance against the uniform law tests whether learning improves on executable support alone, and performance against the empirical-family law tests whether the reference law captures state-dependent preference beyond global family frequency. Neither comparison is interpreted in terms of physical molecular kinetics or synthetic-route likelihood.

\begin{table}[h]
\centering
{\scriptsize\setlength{\tabcolsep}{4pt}
\begin{tabular}{lrrrrrr}
\toprule
Capability cell & $n$ & unif. & emp.\ fam. & \emph{emp+id} & \emph{fam+unif} & $R_\theta$ \\
\midrule
\op{atom\_restate}: element identity change      & 725 & 6.21 & 6.54 & 4.84 & 7.28 & 5.58 \\
\op{atom\_delete}: leaf death                    & 601 & 6.20 & 3.36 & 2.68 & 3.69 & 3.01 \\
\op{atom\_insert}: one-neighbor birth            & 572 & 6.14 & 7.12 & 5.23 & 6.97 & 5.09 \\
\op{bond\_reroute}: single-atom pendant, cyclic  & 464 & 6.24 & 5.26 & 4.07 & 5.48 & 4.29 \\
\op{bond\_reroute}: multi-atom pendant, cyclic   & 216 & 6.32 & 5.56 & 4.38 & 6.14 & 4.96 \\
\op{ring\_system\_restate}: dearomatization$^{\dagger}$ & 202 & 6.25 & 2.84 & 2.97 & 1.72 & 1.85 \\
\op{ring\_system\_restate}: aromatization$^{\dagger}$   & 165 & 6.51 & 2.85 & 2.34 & 1.12 & 0.61 \\
\op{cycle\_insert}: close to monocyclic$^{\dagger}$     & 138 & 6.47 & 7.19 & 3.27 & 4.80 & 0.88 \\
\op{cycle\_attach}: open from monocyclic$^{\dagger}$    & 104 & 6.16 & 5.24 & 5.20 & 3.34 & 3.29 \\
\op{bond\_reorder}: bond-order decrease          &  83 & 6.32 & 4.30 & 4.01 & 3.18 & 2.89 \\
\op{bond\_reorder}: bond-order increase          &  83 & 6.49 & 4.71 & 4.40 & 3.29 & 2.99 \\
\op{cycle\_insert}: close to polycyclic$^{\dagger}$     &  81 & 6.55 & 7.37 & 3.95 & 5.26 & 1.84 \\
\op{cycle\_attach}: open from polycyclic$^{\dagger}$    &  69 & 6.42 & 5.82 & 5.71 & 4.08 & 3.96 \\
\midrule
Overall                                          & 3545 & 6.25 & 5.37 & 4.10 & 5.17 & \textbf{3.90} \\
\bottomrule
\end{tabular}}
\caption{\small Negative log-probability of the realized canonical successor, by capability cell, lower is better. All five laws are evaluated on the same executable support and the same canonical quotient; only the probability assignment differs. Columns \emph{emp+id} and \emph{fam+unif} are the two decomposition arms defined above: empirical family with learned identity, and learned family with uniform identity. Cells below the preregistered minimum size of $20$ are not reported separately; no scored transition was skipped and no partition failed to compile. $^{\dagger}$Supervision for these cells comes entirely from the synthetic-perturbation lane; see the provenance note below.}
\label{tab:rtheta-cells}
\end{table}

The decomposition is lopsided. Learning only which family fires recovers little, improving on corpus family frequencies by $0.20$ nats overall, whereas learning only where within a family the edit lands recovers most of the effect at $4.10$; the full model adds a further $0.20$ beyond that. In the terms used in \cref{sec:exp-rtheta}, the empirical-family-to-full improvement is $1.4616$ nats and the identity-only arm recovers $1.2674$ of it, or $86.7\%$, against $13.3\%$ for the family gate alone. The reference law's advantage is therefore largely an advantage in placing an edit on a particular molecular site rather than in selecting the operator family. The empirical-family control is also not uniformly strong: on both \texttt{cycle\_insert} cells, on \op{atom\_insert} and on \op{atom\_restate} it is worse than assigning equal probability to every distinct successor, so on birth and insertion families a corpus frequency prior misallocates mass relative to ignorance. On four cells the empirical-family/learned-identity ablation assigns higher probability to the realized successor than the full model, indicating that the learned family gate does not uniformly improve over corpus-level family frequencies and that the identity model carries the result on those families.

Aggregating by supervision provenance sharpens this. On the $2{,}744$ transitions from families with data-backed supervision, the reference law retains a $1.06$-nat advantage over empirical family frequencies ($5.55$ against $4.49$), so the main-text claim survives the obvious confound. On that subset, however, the empirical-family/learned-identity ablation reaches $4.25$ nats and therefore assigns higher probability to the realized successor than the full model does, and the learned-family/uniform-identity arm at $5.79$ is worse than corpus family frequency. The learned family gate is thus not uniformly beneficial, and on data-backed families the identity model carries the result in spite of it. The complementary synthetic-only subset comprises $759$ transitions, where the reference law reaches $1.79$ nats against $4.72$ for empirical family frequencies; the aggregate improvement reported in \cref{sec:exp-rtheta} is correspondingly larger than the data-backed figure. The partition is computed from the emitted family-lane provenance artifact rather than from a hand-specified family list, and the restricted aggregates are reproduced from it.

The cells in which the reference law performs best are those whose supervision is synthetic, as recorded in \cref{app:data}. On the ring and cycle families the model recovers a perturbation of a kind it was trained to invert, and the aggregate improvement should not be read as evidence about molecular kinetics. Resolving instead by the number of available successors gives $4.86$ nats for states with no alternative in the recorded band and $5.51$, $5.44$ and $5.51$ for successor counts of $1$--$4$, $5$--$24$ and $25$ or more under the empirical-family law, against $4.34$, $4.03$, $4.00$ and $4.21$ for the learned-identity arm, so the learned advantage is not an artifact of nearly deterministic rows.

\subsection{Future-reachability evaluation}
\label{app:exp-lookahead}

The panel was sealed under preregistration and opened once. Sixty-seven held-out source--target pairs were constructed whose target is reachable under the frozen executable process within the declared budget; two were excluded from the primary analysis under a committed amendment, giving $65$ pairs, and the full $67$ are retained as a sensitivity analysis. All arms share the reference law, the legal edit sets, the target molecule and the edit budget. A myopic policy selects the locally preferred edit; verified lookahead evaluates candidate successors by whether the target remains recoverable within the remaining budget; three intermediate prioritizers restrict which successors are expanded before lookahead is applied. Continuation work is reported as a fraction of the candidates evaluated by all-candidate lookahead, so recovery and cost are read together. The unit is the source--target pair and intervals are paired bootstrap.

\begin{table}[h]
\centering
{\small\setlength{\tabcolsep}{4pt}
\begin{tabular}{lrrrrr}
\toprule
Arm & Rec. & Rate & Cont./pair & Frac.\ univ. & Gain ret. \\
\midrule
Greedy local decision            & 26/65 & $40.0\%$ & $0.0$  & $0\%$    & --- \\
$R_\theta$ top-1 + lookahead     & 33/65 & $50.8\%$ & $7.3$  & $34.7\%$ & $0.50$ \\
Similarity top-1 + lookahead     & 34/65 & $52.3\%$ & $7.3$  & $34.7\%$ & $0.57$ \\
Similarity top-2 + lookahead     & 37/65 & $56.9\%$ & $10.8$ & $52.1\%$ & $0.79$ \\
Goal-aware $h_\phi$ top-1 + lookahead & 40/65 & $61.5\%$ & $7.2$ & $34.6\%$ & $1.00$ \\
All-candidate lookahead          & 40/65 & $61.5\%$ & $20.5$ & $100\%$  & $1.00$ \\
\bottomrule
\end{tabular}}
\caption{\small Sealed panel, $65$ pairs. \emph{Rec.}: targets recovered. \emph{Cont./pair}: candidate continuations actually evaluated per pair. \emph{Frac.\ univ.}: that count as a share of what all-candidate lookahead evaluates. \emph{Gain ret.}: the fraction of all-candidate lookahead's $14$ rescues the arm also recovers.}
\label{tab:lookahead-full}
\end{table}

Of the $65$ pairs, $14$ are recovered by all-candidate lookahead and not by greedy, and none by greedy and not by all-candidate lookahead. That one-sidedness is a property of the decision rule rather than a finding: the planner overrides greedy only on strict improvement in remaining-budget value and ties keep the greedy action, so an arm cannot lose a target greedy would have reached; across the panel $215$ tied candidates kept the greedy action and no arm regresses against greedy. A discordance test would therefore report a foregone conclusion, and the reported quantities are instead the paired difference in recovery rate, $21.5$ percentage points $[12.3,33.5]$, and the share of greedy's $39$ failures that lookahead rescues, $35.9\%$ $[21.2,52.8]$.

Equal recovery counts do not imply equal rescue sets. The goal-aware prioritizer and all-candidate lookahead both recover $40$ targets but not the same $40$: the prioritizer rescues one pair all-candidate lookahead misses and misses one pair it rescues, overlapping on $13$ of $14$. The similarity and reference-probability arms are strictly nested inside all-candidate lookahead's rescue set, missing six, three and seven of it respectively. Resolving by verified path length, over pairs reachable in four edits ($n=21$) the arms recover $7$, $12$ and $13$ for greedy, all-candidate lookahead and the goal-aware prioritizer; at five edits ($n=34$), $14$, $20$ and $19$; at six edits ($n=10$), $5$, $8$ and $8$. The original hypothesis that the advantage would widen with path length was tested and not supported, and horizon is reported as an observation. Including both excluded pairs, greedy recovers $26/67$ and all-candidate lookahead $41/67$, a headroom of $15$ rather than $14$, and no ordering changes. The panel consumed $1{,}347$ kernel calls.

\subsection{Source-conditioned molecular editing}
\label{app:exp-jin}

The protocol is adopted from the comparator and is given in \cref{app:bench-qed}. Each \method{} candidate is the output of one controlled trajectory terminating prospectively at the first realized state inside the goal region, and intermediate molecules are not retrospectively collected. The reference law is unchanged from \cref{app:exp-rtheta}, and the future-value controller is trained and frozen on data disjoint from the official sources before the cohort is opened. A size-fixed internal ablation removes every successor that changes the current heavy-atom count while holding controller, reference law, sources, seeds and candidate allowance fixed. The unit is the source molecule; the candidate curve is descriptive and carries no significance claim, and it is drawn only over the $k$ actually measured, which is $k=1,\dots,12$ on the prospective panel and $k=1,\dots,20$ once the official cohort is opened. The official cohort is opened once and whatever it returns is banked, including a disappointing result, with no hyperparameter re-chosen afterwards.

A receding one-step terminal tilt of the frozen reference law by the immediate property value of each successor was preregistered as the inference policy and refuted on the disjoint development panel, with $0$ successes in $320$ trajectories, before any official-cohort outcome existed. It is retained as a developmental negative control. What this licenses is narrow: the finite-horizon controller is a pre-outcome amendment consistent with the paper's construction, motivated by a developmental failure of a cheap one-step approximation. It is not a restoration of an originally frozen policy, and compute cost was not the only consideration separating the two.

% ===== GAP: NOT YET RUN =====  [G1] official 800-source cohort
% TO CLOSE: run the official cohort once at 20 candidates per source and report
% the success curve over k=1..20 alongside the size-fixed ablation.
The official-cohort success rate, its candidate curve and the size-fixed ablation are \XXX{}.

\subsection{Multi-objective molecular optimization}
\label{app:exp-pareto}

Preference-conditioned optimization over JNK3 and GSK3$\beta$ follows the comparator's setting, detailed in \cref{app:bench-moo}. A common objective-blind initialization bank of $100$ molecules is fixed before either method is evaluated; both receive a matched supervision budget of $10{,}000$ property labels, the same five preference vectors and the same $5{,}000$-evaluation optimization budget. Average Pareto score is primary, with hypervolume, novelty and top-$K$ diversity guarding against scalar gains obtained by collapsing onto a narrow molecule set. The unit is the run, replicated over the frozen bank. Only the immediate-versus-future-aware contrast isolates the contribution of molecular-process control, because \method{} and the comparator do not share a property predictor.

Before training a multi-objective future-value model, we measured how well its Monte Carlo target can be estimated. On $24$ held-out states with $16$ candidate successors each and $32$ reference continuations per candidate, scored by the frozen property model with no oracle calls, the rank correlation between immediate desirability and true future value has per-state median $0.518$ $[0.435,0.700]$ at four remaining edits and $0.318$ $[0.181,0.421]$ at eight, with intervals resampled over states. The two criteria select different successors on $68.3\%$ $[55.8,80.8]$ and $84.2\%$ $[73.3,93.3]$ of state--preference decisions respectively, and following future value rather than immediate desirability improves realized continuation value by a median $12.2\%$ $[1.6,19.8]$ at four edits and $19.0\%$ $[13.1,25.9]$ at eight.

Target reliability itself degrades with horizon. Estimated by split-half agreement between disjoint sets of continuations, a $32$-rollout target has reliability approximately $0.66$ at four remaining edits and $0.39$ at eight, implying by the Spearman--Brown relation that roughly $65$ and $200$ continuations respectively would be required to reach a reliability of $0.80$. Because supervision quality is horizon-dependent while the measured benefit is not, the supported horizon is fixed from target reliability and computational feasibility rather than from optimization performance, as recorded in \cref{app:models}.

\begin{table}[h]
\centering
\caption{\small\textbf{Two-objective molecular optimization at a matched oracle-label budget.} The estimand is the average score of the top-100 molecules (APS) with novelty and top-$K$ diversity; the independent statistical unit is the run, replicated over the frozen initialization bank. The comparator's published row is \emph{reported context} under an unspecified initialization and an unstated hypervolume reference point, and is never the quantity a \method{} number is described as beating; the matched rerun is the row that adjudicates. Every \method{} cell and every matched-rerun cell is unmeasured.}
\label{tab:multiobj}
\footnotesize
\begin{tabularx}{\textwidth}{@{}lY r r r@{}}
\toprule
Arm & Initialization & APS & Novelty & Diversity\\
\midrule
Comparator, as published \emph{(reported context)} & unspecified & $0.841$ & $100\%$ & $0.768$\\
% source: docs/AMENDMENT_INVERSIONGNN_2OBJ.md ("reported: APS 0.841, novelty 100%, diversity 0.768, HV 0.763 +/- 0.031")
Comparator, matched rerun & frozen common bank & \XXX{} & \XXX{} & \XXX{}\\
\method{}, future-aware \eqref{eq:central} & frozen common bank & \XXX{} & \XXX{} & \XXX{}\\
\method{}, immediate & frozen common bank & \XXX{} & \XXX{} & \XXX{}\\
Four-objective escalation & frozen common bank & \XXX{} & \XXX{} & \XXX{}\\
\bottomrule
\end{tabularx}
\end{table}

% ===== GAP: NOT YET RUN =====  [G3] matched two-objective comparison
% TO CLOSE: fix the reliability criterion and admitted horizon, train the
% multi-objective head on admitted horizons only, then run both COMPOSE arms
% and the matched comparator rerun from the frozen bank at the common budget.

\subsection{Mid-trajectory retargeting}
\label{app:exp-retarget}

We evaluate retargeting on $40$ source-disjoint held-out molecules under potency-first and developability-first histories. Each prefix contains three committed edits within a six-edit horizon and is fixed before the eventual joint objective $B$ is constructed, so every post-switch comparison begins from an identical realized prefix with the same remaining budget. Margins are normalized by the held-in interquartile range and the reported objective is the worst margin of the conjunction, with zero denoting satisfaction. At the switch the median worst margin is $-1.115$ following potency-first prefixes and $-1.810$ following developability-first prefixes, so neither history already satisfies $B$.

\begin{table}[h]
\centering
{\small\setlength{\tabcolsep}{4pt}
\begin{tabular}{llrrrc}
\toprule
Comparison & History & Mean & Median & $95\%$ CI & W/L \\
\midrule
Responsiveness & potency & $+0.748$ & $+0.794$ & $[+0.608,+0.891]$ & 36/0 \\
 & developability & $+1.462$ & $+1.357$ & $[+1.240,+1.691]$ & 40/0 \\[2pt]
Greedy-matched history & potency & $+0.367$ & $+0.338$ & $[+0.251,+0.482]$ & 37/3 \\
 & developability & $+0.211$ & $+0.147$ & $[+0.065,+0.361]$ & 26/14 \\[2pt]
\emph{Value of history} & potency & $+0.302$ & $+0.243$ & $[+0.186,+0.413]$ & 35/5 \\
\emph{(primary)} & developability & $+0.176$ & $+0.127$ & $[+0.069,+0.283]$ & 29/11 \\[2pt]
Price of surprise & potency & $-0.252$ & $-0.123$ & $[-0.364,-0.155]$ & 6/32 \\
 & developability & $-0.428$ & $-0.291$ & $[-0.580,-0.288]$ & 4/36 \\
\bottomrule
\end{tabular}}
\caption{\small Held-out retargeting panel, $40$ source-disjoint sources per history. Rows compare, in order: retargeting against continuing the old goal; greedy retargeting against greedy restart from the source; verified continuation from $x_3$ against verified restart from $x_0$ (the primary estimand); and verified continuation against a clairvoyant controller given $B$ from the first edit. Entries are paired post-switch differences in the registered objective; positive values favor the first strategy named. W/L counts sources on which that strategy is better or worse; ties omitted. Intervals are source-level bootstrap.}
\label{tab:retarget-full}
\end{table}

Retargeting materially redirects both histories, and under the primary verified comparison continuing from the realized state remains preferable to restarting from the source. The mean advantage is $+0.302$ $[0.186,0.413]$ following potency-first prefixes and $+0.176$ $[0.069,0.283]$ following developability-first prefixes, with wins on $35/40$ and $29/40$ sources respectively. The clairvoyant controller remains better, giving the negative price-of-surprise effects shown in the table. Arm-level medians follow the same ordering under both histories: continuation from the realized molecule lies above restart and below clairvoyance, at $-0.016$ and $+0.077$ for greedy and verified retargeting against $-0.532$ and $-0.428$ for the two restart arms and $+0.389$ for clairvoyance under potency-first prefixes, and correspondingly under developability-first prefixes.

The effect decreases from development to held-out confirmation, with the primary differences shrinking from $+0.444$ to $+0.302$ and from $+0.290$ to $+0.176$. The confirmation panel contains genuine losses in both histories rather than a deterministic ordering, which is the property that distinguishes the effect from a structural guarantee. We do not use verified-versus-greedy policy improvement as an estimand because its direction follows from the decision rule itself; prefix revisiting is negligible, occurring for at most one source per arm and history, so the effect is not an artifact of trajectories folding back onto their own prefix.

\subsection{Pathwise constraints}
\label{app:exp-pathwise}

The pathwise experiment remains developmental: the cLogP corridor was selected following an earlier feasibility screen, and a separately frozen held-out confirmation has therefore been specified but not opened. Endpoint-only control constrains only the returned molecule; pathwise control removes any successor outside the corridor before the transition is sampled; a terminal-penalty arm changes the objective but leaves the transition support unchanged. All arms share the reference law, objective, horizon, sources, seeds and evaluation budget, and the unit is the source molecule with source-clustered intervals.

On $24$ eligible sources, endpoint-only trajectories return to the corridor after leaving it in $16/24=0.667$ $[0.458,0.833]$ cases under greedy control and $14/24=0.583$ $[0.375,0.792]$ under verified control; the conditional fractions are $16/24$ and $14/23$, so the small-denominator pathology the unconditional estimand guards against did not arise. These excursions are not confined to the boundary: their median depth is $0.69$ cLogP units and their median duration is four of the six committed edits. An earlier protected-motif experiment produced zero such excursions, confirming that the estimand need not be positive simply because intermediate states are observed.

Restricting the support to the corridor remains viable on this panel. The median source retains $79.8\%$ of its legal successors, only one source reaches the prespecified support-tight threshold, no evaluated state has an empty constrained edit set, and one trajectory terminates. Excluding the support-tight source leaves prevalence at $0.652$ and $0.565$ and terminal costs at $-0.121$ and $-0.041$, so the conclusions are unchanged.

The developmental panel does not resolve the terminal-utility cost of imposing the constraint. The paired change in normalized DRD2 utility is $-0.105$ $[-0.377,0.099]$ under greedy parity, with $5$ wins, $8$ losses and $11$ ties, and $-0.023$ $[-0.182,0.134]$ under verified parity, with $9$ wins, $10$ losses and $5$ ties. Both medians are exactly zero and both intervals span zero, which does not establish equivalence: the greedy interval still permits a loss of $0.377$ held-in interquartile-range units. Zero violations under support restriction is an integrity consequence of masking rather than an empirical performance statistic, and the apparent advantage of future-aware control under the mask, $+0.661$ $[0.428,0.944]$, is a ceiling rather than a null, since greedy pathwise control completed all $24$ sources and left no binary headroom.

\begin{table}[h]
\centering
{\small
\begin{tabular}{lrrr}
\toprule
Constraint handling & Complete-path feasible & Wasted path-infeasible evaluations & Final DRD2 utility \\
\midrule
Endpoint rejection & \XXX{} & \XXX{} & \XXX{} \\
Terminal penalty & \XXX{} & \XXX{} & \XXX{} \\
Support restriction & $100\%$ by construction & $0\%$ path-infeasible & \XXX{} \\
\bottomrule
\end{tabular}}
\caption{\small Three-arm comparison for \cref{sec:exp-pathwise}. The $100\%$ complete-path feasibility and $0\%$ wasted path-infeasible evaluations of the support-restriction arm are consequences of the construction and are \emph{not} empirical superiority statistics: successors leaving the corridor are absent from the controlled kernel, so no such state can be committed. The informative comparisons are the first two rows and the final-utility column.}
\label{tab:pathwise}
\end{table}

The held-out confirmation fixes the corridor at $C=[2.3689,4.4522]$ with the same six-edit horizon and endpoint-only semantics, at $n=48$ sources. It passes only if both criteria clear at full $\alpha$ under intersection-union: the lower bound of the source-clustered interval on hidden-path prevalence must exceed $1/3$, a threshold fixed before the development run, and the lower bound on the mean paired change in terminal utility must exceed $-0.25$ held-in interquartile-range units, with $-0.20$ as a prespecified sensitivity analysis. The panel is sized on the joint criterion, giving power $0.976$ at the $0.25$ margin and $0.874$ at $0.20$. No held-out source has been opened.

A value head retrained to be aware of a longer horizon was built, evaluated and rejected on four separate measurements, and is not the head used anywhere in \cref{sec:experiments}. The consequence for the main text is a claim boundary: any gain from a longer horizon comes from giving the frozen process more edit opportunities, not from a better budget-conditioned value model.

\section{Benchmark and baseline details}
\label{app:alignment}

\subsection{{QED} editing benchmark}
\label{app:bench-qed}

The protocol is adopted verbatim from the comparator: the $800$ test molecules with $\mathrm{QED}\in[0.7,0.8]$ as sources, $20$ candidates per source, a Tanimoto similarity constraint of $0.4$ to the source, and success defined as at least one returned candidate reaching $\mathrm{QED}\in[0.9,1.0]$ while satisfying that constraint. Source selection, fingerprints, property calculations, candidate count and aggregation must be identical for the comparison to stand, and where any of these could not be matched it is stated rather than assumed.

\paragraph{Output contract.} Each source receives exactly $K$ candidate slots. A trajectory that reaches the target region returns the first committed molecule satisfying the benchmark criterion. If no sampled continuation reaches the region, the slot returns the unchanged canonical source solely to preserve the fixed-$K$ benchmark interface; this output necessarily fails because benchmark sources begin below the target QED threshold. No failed slot is resampled or replaced. This contract is frozen before the official cohort is opened and is not revised afterwards.

Every non-source returned molecule is a benchmark success: all $371$ non-extinct trajectories return a molecule in the target region, while all $1{,}165$ extinct slots return the unchanged source and are guaranteed failures. A smaller candidate allowance can disadvantage both success and the diversity metric, because a source with only one successful distinct return contributes zero diversity. Separately, no returned success is identical to its source, so the count is unchanged under the stricter rule of the released success script, which requires similarity strictly below one.

Prior-method values are taken as reported in the comparator's paper and are not re-derived. The latent-variable and reinforcement-learning rows arrive as inherited benchmark context rather than as chosen comparators, and no claim in this paper turns on any of them individually; the insertion/deletion graph-diffusion result at $45.1\%$ is the modern row that matters. Higher rates have been reported on this task outside the comparator's table, and \cref{app:qed-historical} records them so that \cref{tab:qed} is not read as the complete published record.

Diversity is computed with the released evaluation script of \citet{jin2019vjtnn} rather than reimplemented from its prose description, and we audited that script rather than infer its behaviour. It admits a candidate only if it satisfies both the similarity constraint and the property target, deduplicates the admitted molecules, excludes a source with no admitted molecule from the average, and assigns a diversity of exactly zero to a source with only one. Similarity is a Morgan fingerprint of radius two on $2048$ bits with chirality disabled. These choices are not cosmetic. On our panel $19$ of the $70$ solved sources return a single distinct molecule and therefore enter the average as zero, which is the whole difference between the reported $0.393$ and the $0.539$ that averaging over only the $51$ multi-molecule sources would give. The reported value is $0.393$ with a bootstrap interval of $[0.331,0.452]$ over sources. Among solved sources \method{} returns on average $4.34$ distinct successful molecules of the twelve permitted, and no successful candidate is identical to its source, so the success count is unchanged under the stricter rule of the released script, which requires similarity strictly below one. We recomputed drug-likeness and similarity from the returned structures with those same functions rather than trusting the values recorded during sampling; the two agree to machine precision and the admission decision agrees on all $1{,}536$ candidates. All returned strings are canonical, so the string-level deduplication of the released script coincides with deduplication by molecule, and no pair of distinct returned molecules shares a fingerprint under the achiral setting.

One limitation belongs to the comparator rather than to us. \citet{ninniri2025griddd} do not state how the diversity column of their main property-optimization table is computed. Their baseline diversity values are identical to those reported in the graph-to-graph translation line, which places the inherited rows on the scale defined above, but the method's own entry is undocumented; a rule matching the one used here is stated only in their appendix, for a separately adapted protocol with a different number of optimization rounds. We therefore report our value as computed under a named and audited definition, and we do not assert that it and the published entry were produced by the same procedure.

Because every \method{} trajectory terminates prospectively on entering the goal region, a run either returns a molecule satisfying both constraints or goes extinct, and the set the diversity script admits is exactly the set of successful returns. The admitted set is therefore constructed the same way as in the published rows rather than merely being similar to it, and the residual differences are the cohort of $128$ against $800$ and the allowance of twelve candidates against twenty. The allowance difference is not neutral for diversity: a smaller allowance yields fewer second successes per source and so more single-molecule sources contributing zero, which depresses the reported value relative to what the same system would produce at twenty candidates.

The benchmark constrains returned molecules and says nothing about internal successor evaluations, and competing methods use their native inference procedures. The comparison therefore measures task performance under a fixed output allowance, no efficiency claim is made, and inference ledgers are reported separately. The entry currently printed is the $128$-source prospective validation panel at twelve of the twenty permitted candidates: it is not the official cohort, it uses fewer outputs than the benchmark allows, and it is not presented as a matched row.

\subsection{Published record on the {QED} editing task}
\label{app:qed-historical}

\Cref{tab:qed} reproduces the comparator set assembled by GrIDDD. Stronger values have been reported on the same task by the graph-to-graph translation line, and we record them in \cref{tab:qed-historical} rather than leave the main table to be read as the full published record. We checked the primary sources rather than inheriting the numbers second-hand. The translation results use the same $800$ source molecules, the same similarity threshold $\delta=0.4$, and the same allowance of $K=20$ decoded candidates per source \citep{jin2019vjtnn}; the rate frequently quoted as $60.6\%$ is VJTNN+GAN, whereas VJTNN itself is reported at $59.9\%$. \citet{jin2020hierg2g} evaluate on the test sets provided by \citet{jin2019vjtnn} under the same range and similarity constraint with $K=20$, and reproduce the $59.9\%$ row, which is what places their $76.9\%$ on the same cohort.

\begin{table}[!ht]
\centering
{\small
\begin{tabular}{llrrr}
\toprule
Method & Reported by & $K$ & QED success ($\uparrow$) & Diversity ($\uparrow$) \\
\midrule
MMPA      & \citet{jin2019vjtnn}  & 20 & 32.9\% & 0.236 \\
VSeq2Seq  & \citet{jin2019vjtnn}  & 20 & 58.5\% & 0.331 \\
VJTNN     & \citet{jin2019vjtnn}  & 20 & 59.9\% & 0.373 \\
VJTNN+GAN & \citet{jin2019vjtnn}  & 20 & 60.6\% & 0.376 \\
AtomG2G   & \citet{jin2020hierg2g} & 20 & 73.6\% & 0.421 \\
HierG2G   & \citet{jin2020hierg2g} & 20 & 76.9\% & 0.477 \\
\bottomrule
\end{tabular}}
\caption{\small\textbf{Further published results on the same similarity-constrained {QED} editing task.} All rows use the $800$-source cohort of \citet{jin2019vjtnn} with $\delta=0.4$ and $K=20$. These are supervised graph-to-graph translation models trained on parallel source-target pairs curated for the target property, whereas the \method{} reference law is trained once without reference to QED and the property enters only through control. The rows therefore measure different problem settings and we do not present them as an ablation of one another, nor is any \method{} claim stated relative to them.}
\label{tab:qed-historical}
\end{table}

\subsection{Two-objective benchmark}
\label{app:bench-moo}

The comparator reports an average Pareto score of $0.841$, novelty of $100\%$, diversity of $0.768$ and hypervolume of $0.763\pm0.031$. Three quantities needed to reproduce those numbers are not recoverable from the paper. The initialization is unspecified: an initialization per weight vector is referred to, but no bank, distribution or seed is given, and which molecules a preference-conditioned optimizer starts from changes the difficulty of the benchmark. The hypervolume reference point for the molecular tasks is never stated, so a hypervolume computed by us is not on the comparator's scale. The released implementation appears to evaluate a single preference vector, against the five described in the paper.

Rather than guess at the missing initialization, we freeze a common, deterministic, objective-blind bank before either method is evaluated and run the comparator ourselves from that bank under a matched label budget, the same five preferences and the same optimization budget. The comparator's published values are reported as context under a different and under-specified initialization; the matched rerun is the row that adjudicates, and the published figure is never the quantity a \method{} number is described as beating.

Average Pareto score, novelty and top-$K$ diversity are primary because they are computed identically from a returned molecule set and require no shared reference geometry. Hypervolume is secondary under a preregistered origin reference point and is not comparable to the published value, whose reference point is unknown; that non-comparability is stated wherever the number appears. Within \method{}, the immediate and future-aware arms share the reference process, legal edit sets, property model, preferences, bank and budget, which is what makes their difference attributable to control. Across methods, no property predictor and no compute accounting is shared, and no cross-method difference is attributed to the search procedure alone.

\section{Software, randomization, and compute}
\label{app:repro}

\subsection{Software and environment}

Every job runs in a containerized Debian image pinned to Python $3.11$ with \texttt{torch} $2.4.0$, \texttt{numpy} $1.26.4$, \texttt{scipy} $1.13.1$ and \texttt{rdkit} $2024.3.5$. Chemistry is pinned to a single RDKit release because canonical molecular identity, ring perception and sanitization all come from it and all three enter the definition of the state space in \cref{app:statespace}.

Only artifacts satisfying the evidence criteria used in the manuscript are cited for scientific claims. Developmental artifacts lacking complete producer or input metadata are retained for audit purposes but excluded from claim-bearing results. Bulk artifacts such as corpora, embeddings, per-run records and trained heads live on a persistent volume rather than in the source repository, while result summaries are versioned alongside the code, so a fresh checkout can re-run every job but cannot read a banked result without the volume. The reproducibility package records the code revision, input identities and result summaries required to regenerate each reported analysis.

\subsection{Randomization and statistics}

Run seeds are derived rather than drawn, by hashing a protocol tag together with the source molecule and the replicate index and taking the leading eight bytes of the digest. In paired comparisons the arm name is deliberately excluded from the seed, so arms sharing a source and replicate index consume identical randomness and a difference between them cannot be a seeding artifact. Where an arm-specific seed is unavoidable for later replicates, the first replicate is still constructed to be identical across arms.

Intervals are bootstrap intervals resampled over the independent unit named in each experiment. That unit is the source molecule in most cases, the source--target pair in the sealed lookahead panel, and the state in the reliability analysis of \cref{app:exp-pareto}. Replicate counts differ by analysis, ranging from $4{,}000$ to $8{,}192$ draws, and each artifact records its own. Where an interval is exact rather than bootstrapped, that is stated with the reason.

\subsection{Compute}

All reported jobs are CPU-only; no result in this paper required a GPU. Sampling containers reserve one core and burst to eight under a six-hour ceiling, and driver containers hold a quarter core under a twelve-hour ceiling. Compute is reported in two units that are never pooled: kernel calls, meaning evaluations of the learned reference law, and container-seconds. The sealed lookahead panel consumed $1{,}347$ kernel calls and roughly $24{,}300$ container-seconds, and retargeting branch points cost on the order of $10^2$ kernel calls and $10^3$ container-seconds each. Oracle calls are counted separately again, because the benchmark budgets in \cref{app:alignment} constrain returned molecules rather than internal evaluations.
